\section{AHSME 1990}
        \begin{enumerate}
        	\item (\textit{AHSME 1990}) If $\dfrac{\frac{x}{4}}{2}=\dfrac{4}{\frac{x}{2}}$, then $x=$


 $\textbf{(A)}\ \pm\frac{1}{2}\qquad\textbf{(B)}\ \pm 1\qquad\textbf{(C)}\ \pm 2\qquad\textbf{(D)}\ \pm 4\qquad\textbf{(E)}\ \pm 8$


 

	\item (\textit{AHSME 1990}) $\left(\frac{1}{4}\right)^{-\frac{1}{4}}=$


 $\textbf{(A) } -16\qquad \textbf{(B) } -\sqrt{2}\qquad \textbf{(C) } -\frac{1}{16}\qquad \textbf{(D) } \frac{1}{256}\qquad \textbf{(E) } \sqrt{2}$


 

	\item (\textit{AHSME 1990}) The consecutive angles of a trapezoid form an arithmetic sequence. If the smallest angle is $75^\circ$, then the largest angle is


 $\textbf{(A) } 95^\circ\qquad \textbf{(B) } 100^\circ\qquad \textbf{(C) } 105^\circ\qquad \textbf{(D) } 110^\circ\qquad \textbf{(E) } 115^\circ$


 

	\item (\textit{AHSME 1990}) 
                    \begin{figure}[H]
                        \centering
                        \includegraphics[width=0.5\textwidth]{Problem/AHSME_I_1990_4_Problem_Figure_1.png}
                    \end{figure}
                    


 Let $ABCD$ be a parallelogram with $\angle{ABC}=120^\circ, AB=16$ and $BC=10.$ Extend $\overline{CD}$ through $D$ to $E$ so that $DE=4.$ If $\overline{BE}$ intersects $\overline{AD}$ at $F$, then $FD$ is closest to


 $\textbf{(A) } 1\qquad \textbf{(B) } 2\qquad \textbf{(C) } 3\qquad \textbf{(D) } 4\qquad \textbf{(E) } 5$


 

	\item (\textit{AHSME 1990}) Which of these numbers is largest?


 $\textbf{(A) } \sqrt{\sqrt[3]{5\cdot 6}}\qquad \textbf{(B) } \sqrt{6\sqrt[3]{5}}\qquad \textbf{(C) } \sqrt{5\sqrt[3]{6}}\qquad \textbf{(D) } \sqrt[3]{5\sqrt{6}}\qquad \textbf{(E) } \sqrt[3]{6\sqrt{5}}$


 

	\item (\textit{AHSME 1990}) Points $A$ and $B$ are $5$ units apart. How many lines in a given plane containing $A$ and $B$ are $2$ units from $A$ and $3$ units from $B$?


 $\textbf{(A) } 0\qquad \textbf{(B) } 1\qquad \textbf{(C) } 2\qquad \textbf{(D) } 3\qquad \textbf{(E) }\text{more than }3$


 

	\item (\textit{AHSME 1990}) A triangle with integral sides has perimeter $8$. The area of the triangle is


 $\textbf{(A) } 2\sqrt{2}\qquad \textbf{(B) } \frac{16}{9}\sqrt{3}\qquad \textbf{(C) } 2\sqrt{3} \qquad \textbf{(D) } 4\qquad \textbf{(E) } 4\sqrt{2}$


 

	\item (\textit{AHSME 1990}) The number of real solutions of the equation

 \[|x-2|+|x-3|=1\]
is


 $\textbf{(A) } 0\qquad \textbf{(B) } 1\qquad \textbf{(C) } 2\qquad \textbf{(D) } 3\qquad \textbf{(E) } \text{more than } 3$


 

	\item (\textit{AHSME 1990}) Each edge of a cube is colored either red or black. Every face of the cube has at least one black edge. The smallest number possible of black edges is


 $\textbf{(A) } 2\qquad \textbf{(B) } 3\qquad \textbf{(C) } 4\qquad \textbf{(D) } 5\qquad \textbf{(E) } 6$


 

	\item (\textit{AHSME 1990}) An $11\times 11\times 11$ wooden cube is formed by gluing together $11^3$ unit cubes. What is the greatest number of unit cubes that can be seen from a single point?


 $\textbf{(A) }  328\qquad \textbf{(B) }  329\qquad \textbf{(C) }  330\qquad \textbf{(D) }  331\qquad \textbf{(E) }  332$


 

	\item (\textit{AHSME 1990}) How many positive integers less than $50$ have an odd number of positive integer divisors?


 $\textbf{(A) } 3\qquad \textbf{(B) } 5\qquad \textbf{(C) } 7\qquad \textbf{(D) } 9\qquad \textbf{(E) } 11$


 

	\item (\textit{AHSME 1990}) Let $f$ be the function defined by $f(x)=ax^2-\sqrt{2}$ for some positive $a$. If $f(f(\sqrt{2}))=-\sqrt{2}$ then $a=$


 $\textbf{(A) } \frac{2-\sqrt{2}}{2}\qquad \textbf{(B) } \frac{1}{2}\qquad \textbf{(C) } 2-\sqrt{2}\qquad \textbf{(D) } \frac{\sqrt{2}}{2}\qquad \textbf{(E) } \frac{2+\sqrt{2}}{2}$


 

	\item (\textit{AHSME 1990}) If the following instructions are carried out by a computer, which value of $X$ will be printed because of instruction $5$?


 <pre>1. START $X$ AT $3$ AND $S$ AT $0$.  
2. INCREASE THE VALUE OF $X$ BY $2$.   
3. INCREASE THE VALUE OF $S$ BY THE VALUE OF $X$. 
4. IF $S$ IS AT LEAST $10000$,   
       THEN GO TO INSTRUCTION $5$;  
       OTHERWISE, GO TO INSTRUCTION $2$.  
       AND PROCEED FROM THERE.  
5. PRINT THE VALUE OF $X$.  
6.STOP. 
</pre>$\textbf{(A) } 19\qquad \textbf{(B) } 21\qquad \textbf{(C) } 23\qquad \textbf{(D) } 199\qquad \textbf{(E) } 201$


 

	\item (\textit{AHSME 1990}) 
                    \begin{figure}[H]
                        \centering
                        \includegraphics[width=0.5\textwidth]{Problem/AHSME_I_1990_14_Problem_Figure_1.png}
                    \end{figure}
                    


 An acute isosceles triangle, $ABC$, is inscribed in a circle. Through $B$ and $C$, tangents to the circle are drawn, meeting at point $D$. If $\angle{ABC}=\angle{ACB}=2\angle{D}$ and $x$ is the radian measure of $\angle{A}$, then $x=$


 $\textbf{(A) } \frac{3\pi}{7}\qquad \textbf{(B) } \frac{4\pi}{9}\qquad \textbf{(C) } \frac{5\pi}{11}\qquad \textbf{(D) } \frac{6\pi}{13}\qquad \textbf{(E) } \frac{7\pi}{15}$


 

	\item (\textit{AHSME 1990}) Four whole numbers, when added three at a time, give the sums $180,197,208$ and $222$. What is the largest of the four numbers?


 $\textbf{(A) } 77\qquad \textbf{(B) } 83\qquad \textbf{(C) } 89\qquad \textbf{(D) } 95\qquad \textbf{(E) cannot be determined from the given information}$


 

	\item (\textit{AHSME 1990}) At one of George Washington's parties, each man shook hands with everyone except his spouse, and no handshakes took place between women. If $13$ married couples attended, how many handshakes were there among these $26$ people?


 $\textbf{(A) } 78\qquad \textbf{(B) } 185\qquad \textbf{(C) } 234\qquad \textbf{(D) } 312\qquad \textbf{(E) } 325$


 

	\item (\textit{AHSME 1990}) How many of the numbers, $100,101,\cdots,999$ have three different digits in increasing order or in decreasing order?


 $\textbf{(A) } 120\qquad \textbf{(B) } 168\qquad \textbf{(C) } 204\qquad \textbf{(D) } 216\qquad \textbf{(E) } 240$


 

	\item (\textit{AHSME 1990}) First $a$ is chosen at random from the set $\{1,2,3,\cdots,99,100\}$, and then $b$ is chosen at random from the same set. The probability that the integer $3^a+7^b$ has units digit $8$ is


 $\textbf{(A) } \frac{1}{16}\qquad \textbf{(B) } \frac{1}{8}\qquad \textbf{(C) } \frac{3}{16}\qquad \textbf{(D) } \frac{1}{5}\qquad \textbf{(E) } \frac{1}{4}$


 

	\item (\textit{AHSME 1990}) For how many integers $N$ between $1$ and $1990$ is the improper fraction $\frac{N^2+7}{N+4}$ $\underline{not}$ in lowest terms?


 $\textbf{(A) } 0\qquad \textbf{(B) } 86\qquad \textbf{(C) } 90\qquad \textbf{(D) } 104\qquad \textbf{(E) } 105$


 

	\item (\textit{AHSME 1990}) 
                    \begin{figure}[H]
                        \centering
                        \includegraphics[width=0.5\textwidth]{Problem/AHSME_I_1990_20_Problem_Figure_1.png}
                    \end{figure}
                    


 In the figure $ABCD$ is a quadrilateral with right angles at $A$ and $C$. Points $E$ and $F$ are on $\overline{AC}$, and $\overline{DE}$ and $\overline{BF}$ are perpendicular to $\overline{AC}$. If $AE=3, DE=5,$ and $CE=7$, then $BF=$


 $\textbf{(A) } 3.6\qquad \textbf{(B) } 4\qquad \textbf{(C) } 4.2\qquad \textbf{(D) } 4.5\qquad \textbf{(E) } 5$


 

	\item (\textit{AHSME 1990}) Consider a pyramid $P-ABCD$ whose base $ABCD$ is square and whose vertex $P$ is equidistant from $A,B,C$ and $D$. If $AB=1$ and $\angle{APB}=2\theta$, then the volume of the pyramid is


 $\textbf{(A) } \frac{\sin(\theta)}{6}\qquad \textbf{(B) } \frac{\cot(\theta)}{6}\qquad \textbf{(C) } \frac{1}{6\sin(\theta)}\qquad \textbf{(D) } \frac{1-\sin(2\theta)}{6}\qquad \textbf{(E) } \frac{\sqrt{\cos(2\theta)}}{6\sin(\theta)}$


 

	\item (\textit{AHSME 1990}) If the six solutions of $x^6=-64$ are written in the form $a+bi$, where $a$ and $b$ are real, then the product of those solutions with $a>0$ is


 $\textbf{(A) } -2\qquad \textbf{(B) } 0\qquad \textbf{(C) } 2i\qquad \textbf{(D) } 4\qquad \textbf{(E) } 16$


 

	\item (\textit{AHSME 1990}) If $x,y>0, \log_y(x)+\log_x(y)=\frac{10}{3} \text{ and } xy=144,\text{ then }\frac{x+y}{2}=$


 $\textbf{(A) } 12\sqrt{2}\qquad \textbf{(B) } 13\sqrt{3}\qquad \textbf{(C) } 24\qquad \textbf{(D) } 30\qquad \textbf{(E) } 36$


 

	\item (\textit{AHSME 1990}) All students at Adams High School and at Baker High School take a certain exam. The average scores for boys, for girls, and for boys and girls combined, at Adams HS and Baker HS are shown in the table, as is the average for boys at the two schools combined. What is the average score for the girls at the two schools combined?

 \[\begin{tabular}[t]{|c|c|c|c|} \multicolumn{4}{c}{Average Scores}\\\hline Category&Adams&Baker&Adams\&Baker\\\hline Boys&71&81&79\\ Girls&76&90&?\\ Boys\&Girls&74&84& \\\hline \end{tabular}\]
$\textbf{(A) } 81\qquad \textbf{(B) } 82\qquad \textbf{(C) } 83\qquad \textbf{(D) } 84\qquad \textbf{(E) } 85$


 

	\item (\textit{AHSME 1990}) Nine congruent spheres are packed inside a unit cube in such a way that one of them has its center at the center of the cube and each of the others is tangent to the center sphere and to three faces of the cube. What is the radius of each sphere?


 $\textbf{(A) } 1-\frac{\sqrt{3}}{2}\qquad \textbf{(B) } \frac{2\sqrt{3}-3}{2}\qquad \textbf{(C) } \frac{\sqrt{2}}{6}\qquad \textbf{(D) } \frac{1}{4}\qquad \textbf{(E) } \frac{\sqrt{3}(2-\sqrt{2})}{4}$


 

	\item (\textit{AHSME 1990}) Ten people form a circle.  Each picks a number and tells it to the two neighbors adjacent to them in the circle.  Then each person computes and announces the average of the numbers of their two neighbors. The figure shows the average announced by each person (\textit{not} the original number the person picked.)

                    \begin{figure}[H]
                        \centering
                        \includegraphics[width=0.5\textwidth]{Problem/AHSME_I_1990_26_Problem_Figure_1.png}
                    \end{figure}
                    
The number picked by the person who announced the average $6$ was


 $\textbf{(A) } 1 \qquad  \textbf{(B) } 5 \qquad  \textbf{(C) } 6 \qquad  \textbf{(D) } 10 \qquad \textbf{(E) }\text{not uniquely determined from the given information}$


 

	\item (\textit{AHSME 1990}) Which of these triples could $\underline{not}$ be the lengths of the three altitudes of a triangle?


 $\textbf{(A) } 1,\sqrt{3},2\qquad \textbf{(B) } 3,4,5\qquad \textbf{(C) } 5,12,13\qquad \textbf{(D) } 7,8,\sqrt{113}\qquad \textbf{(E) } 8,15,17$


 

	\item (\textit{AHSME 1990}) A quadrilateral that has consecutive sides of lengths $70,90,130$ and $110$ is inscribed in a circle and also has a circle inscribed in it. The point of tangency of the inscribed circle to the side of length 130 divides that side into segments of length $x$ and $y$. Find $|x-y|$.


 $\textbf{(A) } 12\qquad \textbf{(B) } 13\qquad \textbf{(C) } 14\qquad \textbf{(D) } 15\qquad \textbf{(E) } 16$


 

	\item (\textit{AHSME 1990}) A subset of the integers $1,2,\cdots,100$ has the property that none of its members is 3 times another. What is the largest number of members such a subset can have?


 $\textbf{(A) } 50\qquad \textbf{(B) } 66\qquad \textbf{(C) } 67\qquad \textbf{(D) } 76\qquad \textbf{(E) } 78$


 

	\item (\textit{AHSME 1990}) If $R_n=\frac{1}{2}(a^n+b^n)$  where $a=3+2\sqrt{2}$ and $b=3-2\sqrt{2}$, and $n=0,1,2,\cdots,$ then $R_{12345}$ is an integer.  Its units digit is


 $\textbf{(A) } 1\qquad \textbf{(B) } 3\qquad \textbf{(C) } 5\qquad \textbf{(D) } 7\qquad \textbf{(E) } 9$


 

\end{enumerate}